\section{IMPLEMENTATION NOTES}
\label{app:implementation}

This appendix records the operational settings needed to reproduce the
experiments. None of it is a contribution of the paper; it is collected here so
that Section~\ref{sec:methodology} can stay on the decisions that affect output
quality.

\subsection{Model Gateway}
\label{app:gateway}

Every model call, including targeted repair calls and the vision calls of
Section~\ref{subsubsec:special}, is issued through one gateway. It fixes two
things: the degree of parallelism, and the response to provider-side failures.
Concurrency is bounded by an asynchronous semaphore, held at 8 in-flight calls in
every experiment so that model comparisons are made under identical conditions.

Responses are classified into three groups. \textbf{Success} (HTTP 2xx) returns
immediately. \textbf{Transient failure} (429, or $\ge$~500) is retried: the
gateway waits for the interval given by the \texttt{Retry-After} header when
present (capped at 60~s), otherwise it applies exponential backoff with
jitter~\citep{exponential-backoff,aws-backoff-jitter},
\[
    \text{delay} = \min(30,\ 2^{\text{retry}}) + \mathrm{U}(0,1)\ \text{seconds},
\]
where $\mathrm{U}(0,1)$ is a uniform random term that prevents many workers from
retrying in lockstep; the retry budget is $N = 2$. \textbf{Permanent failure}
(401, or 400 for other reasons) is reported immediately without retry. Status 429
signals that the client has exceeded the provider's request allowance for the
current interval, and \texttt{Retry-After} carries the provider's own instruction
about when to resume, which takes precedence over any local
schedule~\citep{api-rate-limit-design}. The same policy governs text and vision
calls alike.

The gateway additionally handles responses truncated by the output-token limit
(\texttt{finish\_reason = "length"}) by issuing up to two continuation calls and
stitching the partial JSON fragments together. Per-call service time is observed
by a read-only \texttt{httpx} event hook that does not interfere with control
flow, which is the source of the $T_{\text{req}}$ figures of
Section~\ref{subsubsec:eval-mt}.

\subsection{Parser Runtime}
\label{app:parser-runtime}

The parser processes pages in batches to exploit hardware parallelism. Batch
sizes are set per module (layout~32, text detection~32, text recognition~128,
table-cell detection~256). To prevent VRAM growth and out-of-memory failures, all
intermediate image objects are released at the end of each batch and the GPU
cache is cleared (\texttt{gc.collect}, \texttt{torch.cuda.empty\_cache}). Vision
models are instantiated lazily---loaded on first inference and reused
afterwards---to reduce start-up cost.

\subsection{Fixed Pipeline Parameters}

Table~\ref{tab:app-params} lists the values held constant across all five
systems in the translation experiment, so that the model is the only variable.

\begin{table}[H]
\centering
\caption{Pipeline parameters held fixed across systems.}
\label{tab:app-params}
\small
\begin{tabular}{@{}l l@{}}
\toprule
\textbf{Parameter} & \textbf{Value} \\
\midrule
Concurrent in-flight calls   & 8 \\
Retry budget                 & 2 \\
Sampling temperature         & 0.7 \\
Chunk byte threshold $B$     & 3000 \\
Length tolerance $\tau$      & 0.15 \\
Length-check exemption       & source $<$ 20 characters \\
\midrule
Font-size cluster width $\epsilon$ & 2.0~pt \\
Readability floor $f_{\min}$       & 7.0~pt \\
Table-cell size factor             & 0.88 \\
Downward expansion cap             & 80~pt \\
\midrule
Background ring margin $m$   & 6~pt \\
Appearance raster scale $s$  & $2\times$ \\
Text distance threshold $\theta$ & 80 \\
Centre fraction $\phi$       & 0.6 \\
\bottomrule
\end{tabular}
\end{table}
